\begin{grenze}{Integrierte Grenze v1.6: einfach bedeutet nicht quadratisch}
Für jeden der 60 nativen hellen Zweifermionenzustände ergibt ein passendes besetztes Modenpaar $\langle n_i n_j\rangle=1/8$, während ein Gauß-Zustand mit denselben Zweipunktdaten nur $1/64$ liefern würde. Die exakte Differenz $7/64$ ist frisch für alle 60 Kanäle geprüft. Damit ist der betreffende Zustand in den physischen Fermionmoden nicht gaußsch. Eine andere, nichtlineare Randfeldbeschreibung bleibt möglich, benötigt aber ihre konkrete Übersetzungsregel. Eine Majorana-Zweipunktfunktion allein ersetzt weder dieses Vierpunktdatum noch den nativen Wechselwirkungsoperator.
\end{grenze}
